% !TEX root = ../main.tex
\section{Data: DGS10}
We evaluate on the 10-Year Treasury yield (DGS10), obtained from the Federal Reserve Economic Data (FRED) database via the open-source \texttt{fedfred} Python client \cite{sunderfedfred}, sampled daily from February 1, 2000 through March 26, 2026 (9{,}551 daily observations after gap filling; the final 365 form the holdout beginning March 27, 2025).

\subsection{Rationale}
DGS10 is selected for three reasons.
\paragraph{Macro-financial centrality} Long-term Treasury yields encode expectations of inflation, growth, and monetary policy, making them a canonical macroeconomic indicator.
\paragraph{Known regime shifts} DGS10 exhibits well-documented structural breaks \cite{hamilton1989}, including the zero-lower-bound period and the 2020 COVID shock, providing a natural testbed for nonstationary models.
\paragraph{Multi-scale cyclicality} Treasury yields reflect overlapping cycles (business cycles, policy cycles, liquidity cycles), making them suitable for spectral analysis.
These properties align with prior econometric work emphasizing the importance of structural dynamics in macro time series \cite{hamilton1994}.
